\documentclass[11pt,a4paper,sans]{moderncv}

% ModernCV themes
\moderncvcolor{blue}
\moderncvstyle{classic} 

% Optimized page margins to fit 2 pages
\usepackage[top=1.4cm, bottom=1.4cm, left=1.6cm, right=1.6cm]{geometry}
\setlength{\footskip}{20pt} % Reduced footer space
\definecolor{color1}{rgb}{0.22, 0.45, 0.70}

% Package to control list spacing
\usepackage{enumitem}
\setlist[itemize]{nosep, topsep=0pt, leftmargin=1.5em}

\ifxetexorluatex
  \usepackage{fontspec}
  \usepackage{unicode-math}
  \defaultfontfeatures{Ligatures=TeX}
  \setmainfont{Latin Modern Roman}
  \setsansfont{Latin Modern Sans}
  \setmonofont{Latin Modern Mono}
  \setmathfont{Latin Modern Math} 
\else
  \usepackage[T1]{fontenc}
  \usepackage{lmodern}
\fi
\usepackage[english]{babel}
\usepackage{hyperref}
\usepackage{etoolbox}

% Reduce space after the header title
\patchcmd{\makecvtitle}{\vspace*{1.5em}}{}{}{}

% Personal Data
\name{\textcolor{black}{Chidaksh}}{\textcolor{black}{Ravuru}}
\email{chidakshravuru@gmail.com}
\homepage{https://chidaksh.github.io/}
\social[linkedin]{chidaksh}
\social[github]{chidaksh}

% Adjust column width for dates
\setlength{\hintscolumnwidth}{2.6cm}

%----------------------------------------------------------------------------------
%            Content
%----------------------------------------------------------------------------------
\begin{document}
\makecvtitle

\vspace{-28pt} % Pull up content closer to header

\section{Research Interests}
Developing reliable agentic orchestration and multimodal foundation models for clinical decision support - \nobreak{specially} interested in the alignment, evaluation, and cultural competence of LLMs in high-stakes domains.

\vspace{-6pt}
\section{Education}
\cventry{2024-Present}{University of North Carolina at Chapel Hill}{}{MS in Computer Science}{}{Advisor: \href{https://www.ssriva.com/}{Prof. Shashank Srivastava} --- GPA: 4.0/4.0}
\cventry{2020-2024}{Indian Institute of Technology, Dharwad}{}{B.Tech in Computer Science}{}{Advisors: Prof. Prabuchandran K.J. \& Prof. Rajshekhar Bhat --- GPA: 9.24/10}

\vspace{-6pt}
\section{Selected Publications}
\cventry{2026}{\href{https://arxiv.org/pdf/2512.13142}{\textcolor{color1}{Can LLMs Understand What We Cannot Say? Measuring Multilevel Alignment Through Abortion Stigma Across Cognitive, Interpersonal, and Structural Levels}}}{ACM FAccT}{}{}{Anika Sharma, Malvika Mampally, \textbf{Chidaksh Ravuru}, Kandyce Brennan, Neil Gaikwad}
\cventry{2026}{\textcolor{color1}{Unmasking Spurious Correlations to Fix Text Classifiers}}{Ongoing}{}{}{\textbf{Chidaksh Ravuru}, Shashank Srivastava}
\cventry{2024}{\href{https://arxiv.org/pdf/2408.14484}{\textcolor{color1}{Agentic Retrieval-Augmented Generation for Time Series Analysis}}}{KDD}{}{}{\textbf{Chidaksh Ravuru}, Sagar Srinivas Sakhinana, Venkataramana Runkana}
\cventry{2024}{\href{https://arxiv.org/pdf/2308.14659}{\textcolor{color1}{RESTORE: Graph Embedding Assessment Through Reconstruction}}}{Preprint}{}{}{Hong Yung Yip, \textbf{Chidaksh Ravuru}, Neelabha Banerjee, Shashwat Jha, Amit Sheth, Aman Chadha, Amitava Das}

\vspace{-6pt}
\section{Research Experience}
\cventry{Summer 2025}{Applied Scientist Intern}{Pavo AI}{}{}{
  Designed novel evaluation metrics to capture customer personas into agentic systems 
  \begin{itemize}
    \item Curated PersonaBench, a synthetic dataset of 3,000 conversational trajectories, to train and evaluate AI agents for alignment with specific company personas.
    \item Developed novel evaluation metrics for strategic alignment and used them to fine-tune a reward model, improving an internal planning agent’s persona adherence by 10\%.
  \end{itemize}
}

\cventry{Summer 2023}{Machine Learning Research Intern}{Tata Research Data and Development Centre}{}{}{
  Implemented multi-agentic orchestration system for time series analysis.
  \begin{itemize}
    \item Developed an Agentic-RAG model for time series, achieving a 12\% improvement in prediction accuracy on complex temporal forecasting tasks.
    \item Boosted model performance by 18\% in anomaly detection and 15\% in pattern recognition across 400 task-specific benchmarks using finetuned SLMs.
  \end{itemize}
}

\vspace{-6pt}
\section{Research Experience}
\cventry{Summer 2025}{Applied Scientist Intern}{Pavo AI}{}{}{
  Designed novel evaluation metrics to capture customer personas into agentic systems 
  \begin{itemize}
    \item Curated PersonaBench, a synthetic dataset of 3,000 conversational trajectories, to train and evaluate AI agents for alignment with specific company personas.
    \item Developed novel evaluation metrics for strategic alignment and used them to fine-tune a reward model, improving an internal planning agent’s persona adherence by 10\%.
  \end{itemize}
}

% \cventry{Winter 2024}{Computer Vision Research Intern}{Indian Institute of Technology, Delhi}{}{}{
% Managed the data design, collection, and evaluation for facial recognition in videos.
% \begin{itemize}
%     \item Deployed state-of-the-art face recognition models in law enforcement scenarios and achieved \textbf{85\% Top-1 accuracy}.
%     \item Benchmarked a test dataset comprising \textbf{150+ videos of 50+ subjects} with variations in lighting, height and angle, contributing to research in video-face recognition.
% \end{itemize}}

\cventry{Summer 2023}{Machine Learning Research Intern}{Tata Research Data and Development Centre}{}{}{
  Implemented multi-agentic orcheastration system for time series analysis.
  \begin{itemize}
    \item Developed an Agentic-RAG model for time series, achieving a 12\% improvement in prediction accuracy on complex temporal forecasting tasks.
    \item Boosted model performance by 18\% in anomaly detection and 15\% in pattern recognition across 400 task-specific benchmarks using finetuned SLMs.
  \end{itemize}
}
% \cventry{2019--2021}{Predoctoral Young Investigator}{Allen Institute for AI (AI2)}{}{}{
%   Worked on machine learning and NLP in the AllenNLP group.
%   \begin{itemize}
%     \item Developed and maintained research codebases within the AllenNLP framework.
%     \item Collaborated with researchers to prototype and benchmark new NLP architectures.
%     \item Co-authored internal reports and contributed to publications on large-scale NLP models.
%   \end{itemize}
% }

% \cventry{2018--2019}{Research Assistant}{University of Pennsylvania}{Cognitive Computation Group}{}{
%   Developed evaluations for neural coreference resolution with Prof. Dan Roth.
%   \begin{itemize}
%     \item Implemented evaluation scripts and tailored error analyses for coreference systems.
%     \item Collected and annotated challenging examples to probe model weaknesses.
%     \item Helped design experiments comparing neural and feature-based coreference models.
%   \end{itemize}
% }

% \subsection{Engineering}
% \cventry{Summer 2018}{Software Engineering Intern}{Google, Inc.}{Ads Team}{}{Implemented features for automated bidding systems.}
% \cventry{Summer 2016}{Software Development Intern}{Akuna Capital}{}{}{Developed web apps for compliance and trading visualizations.}
\vspace{-5pt}
\newpage
\section{Projects}
\cventry{2025}{OncoAgentMTB: AI-Powered Molecular Tumor Board}{MedGemma Challenge}{}{}{
\begin{itemize}
    \item Designed a three-stage multi-agent pipeline using Google's HAI-DEF ecosystem: MedGemma 1.5 4B for structured pathology report parsing, MedGemma 27B for multimodal clinical history synthesis, and TxGemma 27B for evidence-based therapeutic recommendations; preliminary evaluation across 10 TCGA cohorts yields 79\% somatic variant classification accuracy and F1 of 0.73 on PubMedQA clinical reasoning tasks.
    \item Automated TCGA genomic-clinical record ingestion and cross-modal data mapping, reducing structured case preparation from $\sim$3 hours to under 30 minutes per patient. 
\end{itemize}}

\cventry{2025}{Persona-Creative Evaluation Pipeline for healthcare creatives}{\href{https://github.com/chidaksh/MedCreatives}{\textcolor{blue}{Code}} $|$ \href{https://github.com/chidaksh/MedCreatives/blob/master/Persona\%20Designing.pdf}{\textcolor{blue}{Report}}}{}{}{
\begin{itemize}
    \item Developed a behavioral evaluation pipeline to study population-level responses to medical creatives across 216+ patient and HCP personas.
    \item Implemented element-level attribution using GPT-4o to identify clinical triggers and skepticism drivers in pharmaceutical messaging.
\end{itemize}}

\cventry{2024}{Explainable Brain Tumor Segmentation with U-Nets}{\href{https://drive.google.com/drive/folders/1k5_SycVpegTsB6DyQMM4N9fFyyx-KPTA?usp=sharing}{\textcolor{blue}{Code}} $|$ \href{https://drive.google.com/file/d/1nbmn6CdoRU7MxpGb7stNprMlg6LcH9y4/view?usp=sharing}{\textcolor{blue}{Report}}}{}{}{
\begin{itemize}
    \item Developed BrainSformer, a transformer architecture for 3D volumetric segmentation, achieving a 95.45\% Dice score on tumor core regions within the BraTS 2020 dataset.
    \item Implemented custom Focal and Dice loss functions to handle class imbalance and integrated Grad-CAM++ for clinical explainability.
\end{itemize}}

\cventry{2024}{Intelligent Planner: Dynamic LLM Query Routing}{\href{https://github.com/chidaksh/PrismLLM}{\textcolor{blue}{Code}} $|$ \href{https://github.com/chidaksh/PrismLLM/blob/master/Intelligent_Planner.pdf}{\textcolor{blue}{Report}}}{}{}{
\begin{itemize}
    \item Architected a dynamic query router using Llama-3.1-8B to orchestrate specialized LLMs, improving inference efficiency and system performance on GSM8K by 3\%.
    \item Engineered the full-stack deployment using Django for backend inference management and React for the frontend.
\end{itemize}}

% \cventry{2024}{Automated Commentary Generation for Soccer Highlights}{\href{https://github.com/chidaksh/SoccerCommentary}{\textcolor{blue}{Code}} $|$ \href{https://arxiv.org/pdf/2508.07543}{\textcolor{blue}{Report}}}{}{}{
% \begin{itemize}
%     \item Addressed a $\sim$50\% BERT score degradation in long-form video models on short highlights by optimizing the temporal window, boosting zero-shot performance by 5.1\%.
%     \item Implemented a progressive fine-tuning strategy to adapt models to short-form video, improving BLEU-4 by 322\% and CIDEr by 72\%.
% \end{itemize}}

% Make sure you have \usepackage{xcolor} in your preamble.

% \section{Technical Side Projects}


% \cventry{2024}{PredictiveAgent: NL2PQL for Kumo.ai}{\href{https://github.com/chidaksh/NL2PQL}{\textcolor{blue}{Code}}}{}{}{
% \begin{itemize}
%     \item Built an autonomous agentic system to convert natural language into Predictive Query Language (PQL) for Relational Foundation Models.
%     \item Designed a Query-First schema architecture that automates graph denormalization and direct link construction for target table definitions.
% \end{itemize}}

% \cventry{2024}{Agentic RAG for Literature Survey and manuscript drafting}{}{}{}{
% \begin{itemize}
%     \item Engineered an end-to-end RAG system to automate academic literature reviews, focusing on the automated retrieval, indexing, and summarization of papers.
%     \item Developed specialized agents to write short, methodology-focused summaries, streamlining the drafting process for academic manuscripts.
% \end{itemize}}

% Make sure you have: \usepackage{xcolor}

\vspace{-5pt}
\section{Teaching}
\cventry{Summer 2025}{Summer School Teacher}{UNC School of Medicine}{}{}{Led a lecture series for high school students on Statistics and Machine Learning (Linear Regression, SVMs) in neuroscience and medical imagery.}
\cventry{2025-Present}{Graduate Teaching Assistant}{UNC School of Data Science and Society}{}{}{TA for DATA 110: Intro to Data Science. Conducted lab recitations and taught basic probability, statistics, and Python coding for undergraduates.}
\cventry{2024}{Undergraduate Teaching Assistant}{Indian Institute of Technology Dharwad}{}{}{TA for Introduction to C, Python, and Software Systems Lab; led workshops and delivered guest lectures.}
% \cventry{Fall 2023}{Teaching Assistant}{UC Berkeley}{}{}{TA for Software Systems Lab and Introduction to C and Python Programming courses.}
% \cventry{2022}{Teaching Assistant}{Winter AI/ML School in India}{Ashoka University}{}{Created assignments for NLP/computer vision.}
% \cventry{2016--2019}{Teaching Assistant}{University of Pennsylvania}{Philadelphia}{}{TA for Mathematical Foundations of CS, Probability, and Algorithms.}

\vspace{-5pt}
\section{Achievements}
\cvitem{2024}{Selected as one of the top 1\% of applicants countrywide for Google Research Week, Bangalore, India}
\cvitem{2023}{Secured 4th place in Inter IIT Techfest 2023 in Open Domain Question Answering}
\cvitem{2022}{Selected for the Machine Learning Summer School (MLSS) 2022 in Krakow, Poland, as one of the few fully funded participants.}

\vspace{-5pt}
\section{Skills}
\cvitem{Programming}{Python, C, C++, JAX, Java, Javascript, Bash/Shell Script, SQL, React, NodeJS, TypeScript, Django, \LaTeX}
\cvitem{Libraries}{PyTorch, HuggingFace Transformers, LangChain, LangGraph, LlamaIndex, TensorFlow, PEFT/LoRA, Pandas, NumPy, Scikit-learn}
\cvitem{Software}{Git, Docker, AWS, Azure, Weights \& Biases, PySpark, MySQL, Unix/Linux, Databricks}

% \begin{rSubsection}{}{}{}{}
% \vspace{-10pt}
% \item {\bf Programming Skills:} C, C++, Python, Java, Bash.
% \item {\bf Python Libraries:} numpy, pandas, matplotlib, seaborn, scipy, sklearn, pytorch and tensorflow.
% \item {\bf Scripts and Software Skills:} HTML, CSS, Javascript, NodeJS, React, PHP, MySQL, LaTeX, Git, Docker
% \item {\bf RESTful API:} LAMP (Linux, Apache, MySQL, Php/python) , Django
% \end{rSubsection}

\vspace{-5pt}
\section{Service \& Outreach}
\cvitem{Reviewing}{ACL, EMNLP, ICLR}
\cvitem{2021-2023}{Appointed as Space Data Science Club Secretary. Directed a team of 30 members, driving data science initiatives focused on space research. Led hands‐on workshops on Python and data visualization, training over 50+ students.}

\end{document}